\scanpage{88}

En \unsure{dernier|deuxième} lieu, soit $f(s)$ une fonction définie pour $s>0$, telle que $f(e^t)$ soit convexe, posons $f(s_1,\mathord{-},s_n)=\sum_1^{\unsure{n}} f(s_i)$, c'est une fonction symétrique des \deleted{\unsure{nombres}}\inserted{\unsure{argts}} des $s_i>0$, et telle que $f(e^{t_1},\mathord{-}e^{t_n})$ soit convexe. Par suite, \deleted{\unsure{$F$}} $f(s_1,\unsure{\mathord{-}|\deleted{+}}s_n)$ satisfait aux conditions du corollaire 1, sous l'hypothèse b. Si on suppose de plus $f(s)$ croissante, et définie et continue aussi pour $s\geq0$, alors $f(s_1,\mathord{-}s_n)$ satisfait $=$ aux conditions du \deleted{corollaire 1}\inserted{\unsure{th.|Th.} 1, 2°} sous l'hypothèse a. D'où, \deleted{\unsure{avec la hypothèse}} \inserted{sous les conditions $\unsure{\rho_i>0}$,} \deleted{\unsure{convenables sur} $(\alpha_i),(\beta_i)$, \unsure{les inégalités}} l'inégalité

\[
\begin{array}{@{}l@{\qquad}l@{}}
(13)&\displaystyle \sum_1^n f(\rho_i)\leq\sum_1^n f(\sigma_i)
\end{array}
\]

et sous les conditions de b), compte tenu d'une remarque précédente, la suite d'inégalités

\[
\begin{array}{@{}l@{\qquad}l@{}}
(14)&\displaystyle \sum_1^k f(\rho_i)\leq\sum_1^k f(\sigma_i)\qquad(k=1,\mathord{-}n)
\end{array}
\]

Prenons p.ex, \deleted{\unsure{$p$ positif et}}\inserted{pour $p$ réel quelconque} $f(s)=s^p$, \inserted{$(s>0)$} on a alors $f(e^t)=e^{pt}$, qui est bien une fonction convexe. \deleted{\unsure{$f$ elle}}\inserted{$f(r)$} est définie et \deleted{\unsure{même}} continue pour $s\geq0$, et \deleted{\unsure{une}} fonction croissante de $s$, dès que $p\geq0$. \deleted{\unsure{les deux} \illegible{1}}

{\reversemarginpar\margin{89}}
